\chapter{対象資料と試験回の整理}
対象は平成23〜30年度、令和元〜7年度、令和元年度再試験の16試験回である。
基礎470問、適性240問、応用理学560問を対象とする。その他部門の保存PDFは教材対象に含めない。
問題・正答の根拠は日本技術士会の公式公開資料\cite{official_questions,official_answers}とし、各問題にPDF物理ページを示す。
\section{年度と試験回}
令和元年度の通常試験と再試験を別回として扱う。試験回を分母にする集計は16、年度を分母にする集計は15となる。
同年度に出題があったかを調べる際に、再試験を別年度として加算しない。
\section{原典へ戻る方法}
問題IDは年度、通常／再試験、科目、番号を組み合わせた識別子である。
PDFの物理ページは1から数える。紙面の印刷ページと表紙・空白ページの有無により一致しない場合がある。
\section{公式正答の例外}
複数の選択肢が正答として受理された問題と、全員へ得点が与えられた問題を区別する。
ALLは全選択肢が内容上正しいという意味ではない。訂正された正答は現行公式値を保持し、訂正前の値と注記を別に追跡する。
